## TEMP holding note — RQ1 plot_level vs. spatial_block integrity check, current tier (2026-08-12)

**Supersedes the old-tier-only version** referenced in `documentation/research_questions_overview.md`
(2026-08-08, `terrain_wind_solid`). This reuses the `spatial_block` side already fit for the
temporal check (`rq1_temporalcheck_fit.sh`) — only `plot_level` needed a new fit
(`rq1_plotlevel_check_fit.sh`, 6 jobs, all succeeded).

### Results

| Model | Cohort | plot_level R2 (easy) | spatial_block R2 | Gap |
|---|---|---:|---:|---:|
| DNN | 4survey | 0.8309 | 0.634 | **0.197** |
| PINN | 4survey | 0.5904 | 0.584 | 0.006 |
| PINN_k | 4survey | 0.5913 | 0.588 | 0.003 |
| DNN | 6survey | 0.8444 | 0.722 | **0.122** |
| PINN | 6survey | 0.7278 | 0.731 | -0.003 |
| PINN_k | 6survey | 0.7275 | 0.730 | -0.002 |

### Headline finding

**Confirms and sharpens the old-tier result.** DNN's inflation on the easy split is even larger
here than the old tier showed (0.197 vs. the old 0.136 on 4survey) — a genuinely large gap, DNN's
plot_level R2 (0.83) badly overstates what `spatial_block_kfold` (0.63) actually shows. PINN and
PINN_k, by contrast, show **essentially zero inflation on either cohort** — gaps of ±0.006 or
smaller, indistinguishable from noise. This is now confirmed on the current VIF-screened tier,
not just the old one: PINN's physics constraint appears to make it structurally unable to exploit
whatever leakage the easy split allows, while DNN readily does. This is a real, citable
methodological justification for using `spatial_block_kfold` throughout, independent of and
additional to the temporal-split finding.
